\documentclass[border=5mm]{standalone}
% ==============================================================================
% SETUP.TEX - CONFIGURACIÓN GLOBAL DEL PROYECTO
% ==============================================================================
% Este archivo centraliza todos los paquetes y estilos.
% Debe incluirse en el preámbulo del libro principal y de los archivos standalone.
% \PassOptionsToPackage{gray}{xcolor}
% --- 1. CODIFICACIÓN E IDIOMA ---
% fix-cm habilita las fuentes Computer Modern en tamaños arbitrarios (por debajo
% de 5 pt, entre otros). Debe cargarse antes que fontenc.
\usepackage{fix-cm}
\usepackage[utf8]{inputenc}
\usepackage[T1]{fontenc}
\usepackage[spanish, es-tabla]{babel}  
\usepackage{csquotes} % Recomendado para babel y citas

\usepackage{import}
% mode=image: Usa el PDF pre-compilado, nunca intenta recompilar.
% Debes compilar cada figura manualmente antes de compilar el libro.
\usepackage[mode=image, subpreambles=false]{standalone}

% --- 2. MATEMÁTICAS Y CIENCIAS ---
\usepackage{amsmath, amssymb, amsfonts, mathtools}
\usepackage{enumitem}

% LaTeX trae \DeclareMathSizes{5}{5}{5}{5}, así que a 5 pt (\tiny) los subíndices
% salen del mismo tamaño que la base: en las etiquetas de figuras el M de
% $\omega_M$ queda desproporcionado. Se restablece la proporción habitual.
\DeclareMathSizes{5}{5}{3.7}{3}

% --- 3. DISEÑO Y GEOMETRÍA --- 
% Unificamos los márgenes para todo el documento
\usepackage[left=2.5cm,right=2.5cm,top=3cm,bottom=3cm]{geometry}
\makeatletter
\ifstandalone % Evita que geometry dañe el recorte de las figuras bajándolas 3cm
  \@ifclasswith{standalone}{crop=false}{
    % Es un capítulo/sección: Dejamos que geometry actúe.
  }{
    % Es una figura: Neutralizamos geometry para que no las recorte mal.
    \geometry{pass}
  }
\fi
\makeatother
\usepackage{parskip} % Espaciado entre párrafos sin sangría
\usepackage{float}   % Para usar [H] en figura

% Ajustes de flotantes para evitar espacios verticales exagerados
\renewcommand{\topfraction}{0.95}      % Permite que las figuras ocupen hasta el 95% superior
\renewcommand{\bottomfraction}{0.95}   % Permite que las figuras ocupen hasta el 95% inferior
\renewcommand{\textfraction}{0.05}     % Permite hojas con tan solo un 5% de texto
\renewcommand{\floatpagefraction}{0.8} % Requiere que una página de solo figuras esté 80% llena
\setcounter{topnumber}{4}
\setcounter{bottomnumber}{4}
\setcounter{totalnumber}{10}

% --- 4. GRÁFICOS Y TABLAS ---
\usepackage{graphicx}
\usepackage{booktabs} % Tablas profesionales
\usepackage{caption}  % Control de captions y \captionof
\usepackage{tikz}
\usetikzlibrary{arrows.meta, calc, angles, quotes, babel, backgrounds}
\usepackage{pgfplots}
\usepgfplotslibrary{groupplots}
\usepgfplotslibrary{external}
\usepgfplotslibrary{patchplots}
\pgfplotsset{compat=1.18}
\usepackage[american, straightvoltages]{circuitikz}

% --- 5. COLORES Y CAJAS ---

\usepackage[table]{xcolor}
\usepackage{tcolorbox}

% Definición de colores corporativos/estilo
\definecolor{utnblue}{RGB}{0, 51, 102}
\definecolor{codegreen}{rgb}{0,0.6,0}
\definecolor{codegray}{rgb}{0.5,0.5,0.5}
\definecolor{codepurple}{rgb}{0.58,0,0.82}
\definecolor{backcolour}{rgb}{0.96,0.96,0.96}

% --- 6. ENTORNOS PERSONALIZADOS (CAJAS) ---

% Caja "Resultado" (Azul Corporativo) — Definiciones, fórmulas clave, resultados importantes.
% Uso: \begin{resultado}[Fórmula de Euler] ... \end{resultado}
\newtcolorbox{resultado}[1][]{
    colback=utnblue!5!white,
    colframe=utnblue,
    title=\textbf{#1},
    fonttitle=\bfseries,
    sharp corners,
    boxrule=0.5mm
}

% Caja "Nota" (Gris) — Advertencias, aplicaciones, contexto secundario.
% Uso: \begin{nota}[Cuidado con la calculadora] ... \end{nota}
% Uso: \begin{nota} ... \end{nota}  → título por defecto "Nota"
\newtcolorbox{nota}[1][Nota]{
    colback=codegray!5!white,
    colframe=codegray,
    title=\textbf{#1},
    fonttitle=\bfseries,
    sharp corners,
    boxrule=0.5mm
}

% Caja inline para resaltar un valor y enlazarlo con su otra aparición en una tabla
% o en una cuenta: mismo color, mismo valor.
% Uso: \marcavalor{$4$}  o  \marcavalor[codegreen]{$4$}
\newtcbox{\marcavalor}[1][utnblue]{
    on line,
    boxsep=1pt, left=2pt, right=2pt, top=1pt, bottom=1pt,
    colback=#1!10!white,
    colframe=#1,
    boxrule=0.4pt,
    arc=2pt
}

% --- 7. CÓDIGO FUENTE (LISTINGS) ---
\usepackage{listings}

\lstdefinestyle{miestilo}{
    backgroundcolor=\color{backcolour},   
    commentstyle=\color{codegreen},
    keywordstyle=\color{magenta},
    numberstyle=\tiny\color{codegray},
    stringstyle=\color{codepurple},
    basicstyle=\ttfamily\footnotesize,
    breakatwhitespace=false,         
    breaklines=true,                 
    captionpos=b,                    
    keepspaces=true,                 
    numbers=left,                    
    numbersep=5pt,                  
    showspaces=false,                
    showstringspaces=false,
    showtabs=false,                  
    tabsize=2,
    frame=single,                    
    rulecolor=\color{codegray}
}

\lstset{style=miestilo}
% Soporte para tildes y ñ en bloques de código
\lstset{literate={ñ}{{\~n}}1 {á}{{\'a}}1 {é}{{\'e}}1 {í}{{\'i}}1 {ó}{{\'o}}1 {ú}{{\'u}}1}

% Operadores matemáticos personalizados

% Vector en sentido discreto (lista finita de coordenadas): negrita + flecha.
% Uso: \vect{v}, \vect{e}_k, \vect{0}.
\newcommand{\vect}[1]{\vec{\mathbf{#1}}}

% Fecha de versión y comando para capítulos standalone
\providecommand{\verdate}{\the\year.\ifnum\month<10 0\fi\the\month.\ifnum\day<10 0\fi\the\day}
\newcommand{\chapterversion}{%
    \vspace{-1.5cm}\begin{flushright}\small\textit{\color{codegray}Versión~\verdate}\end{flushright}\vspace{0.5cm}%
}

% --- 8. HIPERVÍNCULOS (DEBE SER EL ÚLTIMO PAQUETE) ---
\usepackage[colorlinks=true, linkcolor=utnblue, urlcolor=utnblue, citecolor=utnblue]{hyperref}

% --- 9. REFERENCIAS CRUZADAS ENTRE CAPÍTULOS STANDALONE ---
% Cuando se compila un capítulo standalone, xr-hyper lee los .aux de los
% otros capítulos (si existen) para resolver \ref{} a labels externos.
% En la compilación del libro completo este bloque no se activa.
\ifstandalone
  \usepackage{xr-hyper}
  \makeatletter
  \newcommand{\xrchap}[1]{%
    \edef\xrc@self{\detokenize{Capitulo#1}}%
    \edef\xrc@job{\jobname}%
    \ifx\xrc@self\xrc@job\else
      \IfFileExists{../#1capitulo/Capitulo#1.aux}%
        {\externaldocument{../#1capitulo/Capitulo#1}}{}%
    \fi
  }
  \makeatother
  \xrchap{01}\xrchap{02}\xrchap{03}\xrchap{04}
  \xrchap{05}\xrchap{06}\xrchap{07}\xrchap{08}
  \xrchap{09}\xrchap{10}\xrchap{11}\xrchap{12}
  \xrchap{13}
\fi
% \PassOptionsToPackage{gray}{xcolor}

% fig_eje_a_circulo
% Dos columnas: a la izquierda las representaciones sobre la recta, apiladas
% ((a) eje omega y (b) eje lambda); a la derecha la gráfica polar (c).
% (a) Eje omega con el espectro muestreado (copias cada ws), ticks ws/2, ws.
% (b) El MISMO dibujo con el eje medido en lambda = omega*Delta t: las copias
%     quedan cada 2pi y la banda base ocupa [-pi, pi] (franja gris). Es la
%     representación habitual del espectro de una secuencia.
% (c) El eje de (b) enrollado sobre el círculo de longitud 2pi.
% El espectro es un lóbulo suave cuyo módulo decae a cero al aumentar la
% frecuencia:  b(u) = (1 + cos(180 u / W))/2,  |u| < W,  W = ws/5.
% En (c) el lóbulo va en polares con el CENTRO como cero: cada dirección es una
% frecuencia lambda y la distancia al centro es |X_s|, de modo que al aumentar
% lambda el módulo tiende a cero. A < R para que el lóbulo quede DENTRO del
% círculo punteado: ese círculo es el eje de lambda, no una curva de nivel de |X_s|. W = ws/5 son 72 grados en lambda, así que el
% módulo ya vale cero bastante antes de lambda = pi (Nyquist), donde no hay
% contenido. SIN relleno en (c): lo que se lee es la curva del módulo.

\begin{document}
\begin{tikzpicture}[>={Latex[length=4pt]}, font=\footnotesize,
    declare function={
        b(\u)  = 0.5*(1+cos(600*\u));   % lóbulo sobre el eje, |u| < 0.3 cm
        bc(\t) = 0.5*(1+cos(2.5*\t));   % el mismo lóbulo en lambda, |t| < 72 grados
    },
    tri/.style={utnblue, thick, fill=utnblue!14},
    tricopy/.style={utnblue!70, thick, fill=utnblue!6},
    ax/.style={->, thin, black},
    marca/.style={font=\footnotesize},
]
    % ---- La tira del espectro muestreado, idéntica en (a) y en (b) ----
    % #1: nombre del eje. ws = 1.5 cm, semiancho del lóbulo 0.3 cm.
    \newcommand{\tiraespectro}[1]{%
        \draw[ax] (-3.75,0) -- (3.75,0) node[right] {#1};
        \foreach \c in {-3.0, -1.5, 1.5, 3.0} {
            \pgfmathsetmacro{\la}{\c-0.3}
            \pgfmathsetmacro{\lb}{\c+0.3}
            \draw[tricopy] plot[domain=\la:\lb, samples=60, variable=\x]
                (\x,{0.62*b(\x-\c)}) -- cycle;
        }
        \draw[tri] plot[domain=-0.3:0.3, samples=60, variable=\x]
            (\x,{0.62*b(\x)}) -- cycle;
        % eje vertical: el módulo del espectro
        \draw[ax] (0,0) -- (0,0.95);
        \node[gray!75, font=\scriptsize, anchor=east] at (-0.06,0.86) {$|X_s|$};
    }

    % ============ (a) eje omega ============
    % scale sube el dibujo sin agrandar la tipografía: el eje gana cuerpo y las
    % etiquetas quedan del mismo tamaño, con más aire entre ellas.
    \begin{scope}[scale=1.2]
        \def\S{1.5}   % ws en cm
        \tiraespectro{$\omega$}
        % ticks omega
        \draw[thin] (0.5*\S,0.05) -- (0.5*\S,-0.05) node[below, font=\scriptsize] {$\tfrac{\omega_s}{2}$};
        \draw[thin] (\S,0.05) -- (\S,-0.05) node[below, font=\scriptsize] {$\omega_s$};
        \draw[thin] (-0.5*\S,0.05) -- (-0.5*\S,-0.05) node[below, font=\scriptsize] {$-\tfrac{\omega_s}{2}$};
        \draw[thin] (-\S,0.05) -- (-\S,-0.05) node[below, font=\scriptsize] {$-\omega_s$};
        % bracket banda base
        \draw[gray!70, thin] (-0.5*\S,1.12) -- (-0.5*\S,1.22) -- (0.5*\S,1.22) -- (0.5*\S,1.12);
        \node[gray!70, font=\scriptsize, anchor=south] at (0,1.20) {banda base};
        \node[marca] at (0.1,1.68) {(a)};
    \end{scope}

    % ============ Flecha: cambio de unidad ============
    \draw[->, thick, gray!70] (-1.05,-0.55) to[bend right=12] (-1.05,-1.30);
    \node[gray!70, font=\footnotesize, anchor=east] at (-1.15,-0.92)
        {$\lambda = \omega\,\Delta t$};

    % ============ (b) el mismo dibujo, eje en lambda ============
    \begin{scope}[shift={(0,-2.30)}, scale=1.2]
        \def\S{1.5}   % 2pi en cm (el mismo paso que ws en (a))
        % la vuelta [-pi, pi] resaltada: es el tramo que se enrolla en (c)
        \fill[gray!12] (-0.5*\S,0) rectangle (0.5*\S,0.80);
        \tiraespectro{$\lambda$}
        % ticks lambda
        \draw[thin] (0.5*\S,0.05) -- (0.5*\S,-0.05) node[below, font=\scriptsize] {$\pi$};
        \draw[thin] (\S,0.05) -- (\S,-0.05) node[below, font=\scriptsize] {$2\pi$};
        \draw[thin] (-0.5*\S,0.05) -- (-0.5*\S,-0.05) node[below, font=\scriptsize] {$-\pi$};
        \draw[thin] (-\S,0.05) -- (-\S,-0.05) node[below, font=\scriptsize] {$-2\pi$};
        \node[marca] at (0.1,1.15) {(b)};
    \end{scope}

    % ============ Flecha enrollar: de la columna izquierda a la derecha ============
    \draw[->, thick, gray!70] (5.00,-1.15) -- (6.30,-1.15);
    \node[gray!70, font=\footnotesize, anchor=south] at (5.65,-1.05) {enrollar};

    % ============ (c) el círculo lambda, a la derecha ============
    \begin{scope}[shift={(9.20,-1.15)}, scale=0.88]
        \def\R{1.30}
        \def\A{1.15}
        % circunferencia: el eje de frecuencias lambda (una vuelta = 2 pi)
        \draw[gray!55, thin, dashed] (0,0) circle (\R);
        \fill[gray!70] (0,0) circle (0.8pt);
        % el espectro enrollado: distancia al centro = |X_s|, cero en el centro
        \draw[utnblue, thick]
            plot[domain=-72:72, samples=160, variable=\t] (\t:{\A*bc(\t)});
        % marcas del circulo
        \fill[black] (0:\R) circle (1.1pt);
        \node[anchor=west, font=\scriptsize] at (0:{\R+0.10}) {$\lambda=0$};
        \fill[black] (180:\R) circle (1.1pt);
        \node[anchor=east, font=\scriptsize] at (180:{\R+0.10}) {$\lambda=\pi$};
        \node[utnblue, font=\scriptsize, anchor=north east, yshift=-4pt]
            at (180:{\R+0.10}) {$(\omega_s/2)$};
        \draw[gray!70, thin] (90:{\R-0.09}) -- (90:{\R+0.09});
        \node[font=\scriptsize, anchor=south] at (90:{\R+0.14}) {$\tfrac{\pi}{2}$};
        \draw[gray!70, thin] (270:{\R-0.09}) -- (270:{\R+0.09});
        \node[font=\scriptsize, anchor=north] at (270:{\R+0.14}) {$-\tfrac{\pi}{2}$};
        % sentido de lambda creciente
        \draw[->, gray!60, thin] (108:{\R+0.16}) arc (108:150:{\R+0.16});
        \node[gray!75, font=\scriptsize, anchor=south east]
            at (146:{\R+0.20}) {$\lambda$};
        \node[marca] at (0,{\R+1.00}) {(c)};
    \end{scope}
\end{tikzpicture}
\end{document}
